\section{A saturating conserved stress and reduction to a flat end}
\label{sec:str-reduction}

Throughout this section, $A_g(S)$ denotes the full intrinsic-boundary
infimum $A_{\min}(S;g)$ of Definition~\ref{def:setup-cut}. All normals at
the original boundary point into the designated exterior. On an end,
$O_j(r^{-a})$ denotes a bound by $Cr^{-a}$ after at most $j$ derivatives
chosen from $r\partial_r$ and smooth angular vector fields. We write
$O_\infty$ when this holds for each fixed derivative order. Bounds in
physical weighted H\"older spaces are stated separately; they do not,
without further argument, assert bounds for all angular derivatives.

\begin{theorem}[The one-sign component inequality]
\label{thm:str-one-sign}
Let $(\Omega,g,K)$ satisfy Definition~\ref{def:setup-data}, and assume
$H+\tr_{TS}K\le0$ on every component of $S$, with the convention of
Equation~\eqref{eq:setup-expansion}. Suppose that $K$
is compactly supported and that one of the two tensors
\[
 K-(\tr_gK)g,\qquad -K+(\tr_gK)g
\]
is positive semidefinite throughout $\Omega$. In each fixed admissible
hyperbolic chart for which the metric fluxes exist,
\begin{equation}
 p_0(g)\geq \sqrt{\frac{A_g(S)}{16\pi}}
       +\frac12\left(\frac{A_g(S)}{4\pi}\right)^{3/2}.
 \label{eq:str-component-bound}
\end{equation}
This assertion does not require the flux covector to be timelike.
If that covector is future timelike, its Lorentz norm satisfies the
same lower bound.
\end{theorem}

The asymptotically flat result used in the proof is
Theorem~\ref{thm:fl-deformation}. Its required interface is the following:
a smooth exterior with one asymptotically Euclidean end, the full-cut
convention above, a strictly negative future boundary expansion, and
\begin{equation}
 \begin{split}
 R_q+(\tr_q B)^2-|B|_q^2
 -2\big|\operatorname{div}_q(B-(\tr_qB)q)\big|_q
     &\geq c(1+s)^{-3-\delta},\\
 q-\delta_{\mathrm{Eucl}}=O_\infty(s^{-1}),\qquad
 R_q&=O_\infty(s^{-3-\delta}),\\
 B=O_\infty(s^{-2}),\qquad \tr_qB&=0
 \quad\hbox{outside a compact set},
 \end{split}
 \label{eq:str-flat-interface}
\end{equation}
where $c>0$ and $0<\delta<1$, has ADM metric energy at least
$\sqrt{A_q(S)/(16\pi)}$. That theorem is proved in this manuscript; it
is not an additional premise.

\subsection{Conformal preparation with preservation of the component mass}

It is useful to fix the constraint normalization
\begin{equation}
 \mathcal H_g(K)=R_g+6+(\tr_gK)^2-|K|_g^2,
 \qquad
 \mathcal J_g(K)=\operatorname{div}_g(K-(\tr_gK)g).
 \label{eq:str-constraint-normalization}
\end{equation}
The dominant energy condition is $\mathcal H_g(K)\geq2|\mathcal J_g(K)|_g$.

\begin{lemma}[Preparation]
\label{lem:str-preparation}
Under the hypotheses of Theorem~\ref{thm:str-one-sign}, there are smooth
data $(g_i,K_i)$ on the same exterior, numbers $c_i>0$, and a fixed
$0<\delta<1$ such that:
\begin{enumerate}[label=(\roman*)]
\item $K_i$ is compactly supported, with the same global stress sign as
      $K$, and its future boundary expansion has the original weak sign;
\item $\mathcal H_{g_i}(K_i)-2|\mathcal J_{g_i}(K_i)|_{g_i}
      \geq c_i(1+r)^{-3-\delta}$, with $r$ smoothly extended over a
      compact region;
\item $g_i=e^{4v_i}b$ on an exact tail, where
\begin{equation}
 v_i=r^{-3}v_{3,i}(\omega)-d_{0,i}r^{-3-\delta}
             +O_\infty(r^{-4}),
 \qquad d_{0,i}>0,
 \label{eq:str-prepared-expansion}
\end{equation}
and $v_{3,i}$ is smooth;
\item for a fixed $\beta>3/2$, $g_i\to g$ and $K_i\to K$ in physical
      weighted $C^{2,\alpha'}_\beta$ and $C^{1,\alpha'}_\beta$ norms,
      respectively, with any fixed sufficiently small $\alpha'>0$;
\item $p_a(g_i)\to p_a(g)$ for all four metric fluxes and
      $A_{g_i}(S)\to A_g(S)$.
\end{enumerate}
When $K=0$, one may take $K_i=0$. A minimal boundary remains minimal,
and the convergence on fixed compact subsets is smooth if the original
metric is smooth there.
\end{lemma}

\begin{proof}
Choose
\begin{equation}
 \frac32<\beta<\min\{\tau,3\},\qquad
 0<\delta<\min\{2\beta-3,1\},\qquad
 \eta_R=R^{\beta-\tau}.
 \label{eq:str-preparation-parameters}
\end{equation}
Take $R$ beyond the support of $K$ and a fixed cutoff
$\chi_R=\chi(r/R)$, equal to one for $r\leq R$ and zero for $r\geq2R$.
Set $q_R=b+\chi_R(g-b)$ on the end and $q_R=g$ elsewhere. For large
$R$ this is positive definite. The curvature cutoff error
\[
 E_R=\frac18\{\chi_R(R_g+6)-(R_{q_R}+6)\}
\]
is supported in $R\leq r\leq2R$ and satisfies
\begin{equation}
 \|E_R\|_{C^{0,\alpha'}_\beta}\leq C\eta_R.
 \label{eq:str-cutoff-source}
\end{equation}
Indeed, the scalar-curvature formula uses two derivatives of the
metric, physical derivatives of $\chi(r/R)$ are uniformly bounded on
this annulus, and every term contains $g-b$ or one of its first two
physical derivatives. The quadratic terms decay faster. The metrics
$q_R$ have uniform bounded-geometry coordinate estimates of the orders
used here.

Choose a nonnegative smooth compactly supported function $D\geq|K|_g$,
with support in the region where $q_R=g$, and a smooth positive $w_0$
equal to $r^{-3-\delta}$ far out. We solve
\begin{equation}
 \begin{cases}
 (-\Delta_{q_R}+3)v
 =D\sqrt{\eta_R^2+|dv|_{q_R}^2}+E_R+\eta_Rw_0,\\
 \partial_\nu v=0\quad\hbox{on }S,\qquad v\longrightarrow0
 \quad\hbox{at infinity}.
 \end{cases}
 \label{eq:str-preparation-equation}
\end{equation}
Here the sign of the normal in the homogeneous Neumann condition is
immaterial.

We give the existence and uniform estimates. On a finite truncation
put $v=0$ on its outer sphere, and multiply the entire right side of
Equation~\eqref{eq:str-preparation-equation} by $\lambda\in[0,1]$.
Write its gradient term as $W\cdot dv+F_D$, where $|W|\leq D$ and
$0\leq F_D\leq D\eta_R$; for example this follows from
\[
 \sqrt{\eta_R^2+|dv|^2}
 =\frac{|dv|^2}{\sqrt{\eta_R^2+|dv|^2}}
   +\frac{\eta_R^2}{\sqrt{\eta_R^2+|dv|^2}}.
\]
The maximum principle for $-\Delta-W\cdot d+3$, with the indicated
mixed boundary conditions, bounds $|v|$ by $C\eta_R$. Outside a fixed
sphere containing $\supp D$, the radial function $r^{-\beta}$ satisfies
\[
 (-\Delta_{q_R}+3)r^{-\beta}
 =\{(3-\beta)(\beta+1)+o(1)\}r^{-\beta}
 \geq c_\beta r^{-\beta}
\]
uniformly in $R$. Comparison using Equation~\eqref{eq:str-cutoff-source}
therefore sharpens the bound to
\begin{equation}
 |v|\leq C\eta_R(1+r)^{-\beta}.
 \label{eq:str-v-height}
\end{equation}
The boundary constants in these comparisons are independent of the
outer truncation.

For completeness, the required derivative bounds follow from scalar
elliptic estimates without any coupled-system assumption. Divide the
equation by $\eta_R$ and write $u=v/\eta_R$. Its gradient dependence is
$D\sqrt{1+|du|^2}$, with uniformly bounded first derivative in $du$.
Local $W^{2,p}$ estimates on physical balls, the corresponding Neumann
or Dirichlet half-ball estimates, and first-derivative interpolation
give a bound for $u$ in $W^{2,p}$ on smaller patches in terms of its
already bounded height and the source $E_R/\eta_R+w_0$. One can see the
absorption directly by applying the estimate to nested cutoffs and
using
\[
 \|du\|_{L^p}\leq\varepsilon\|u\|_{W^{2,p}}
                       +C_\varepsilon\|u\|_{L^p}.
\]
Choose $p>3$ sufficiently large. Sobolev embedding makes $du$ H\"older,
and the scalar Schauder estimate then gives $C^{2,\alpha'}$ control.
On the end the same argument is applied after multiplication by the
local weight, which is comparable on every physical unit patch.
Thus
\begin{equation}
 \|v\|_{C^{2,\alpha'}_\beta}\leq C\eta_R.
 \label{eq:str-v-estimate}
\end{equation}
Lemma~\ref{lem:linear-separated-faces} supplies these scalar estimates
and the finite-domain inverse. Its local reflection argument permits
bounded drift, and its Schauder step is applied to the original
one-sided equation after the gradient becomes H\"older.
The ellipticity, boundary smoothness, source regularity and
gradient-growth hypotheses have just been specified. The linearization of the finite-domain equation is
\[
 -\Delta_{q_R}+3-
 \lambda D\frac{\langle dv,d(\,\cdot\,)\rangle}
                    {\sqrt{\eta_R^2+|dv|^2}}.
\]
Its kernel vanishes by the same maximum principle. Its index is zero
by continuation of the bounded drift from zero. It is therefore an
isomorphism with these boundary conditions. The a priori bounds close
the continuation from $\lambda=0$ to $1$. Exhaustion, the uniform
estimates, and a diagonal subsequence give a smooth solution of
Equation~\eqref{eq:str-preparation-equation}.

On its exact hyperbolic tail, this solution satisfies
$(-\Delta_b+3)v=\eta_Rr^{-3-\delta}$. Rotations commute with the
operator and annihilate this radial source. Applying difference
quotients, comparison by $r^{-\beta'}$ for any fixed $\beta'<3$, and
local elliptic estimates gives the same decay for all fixed angular
derivative orders. The hyperbolic Laplacian is
\[
 \Delta_b=(1+r^2)\partial_r^2+(2/r+3r)\partial_r
                         +r^{-2}\Delta_\sigma.
\]
Consequently, for $\beta'$ sufficiently close to $3$,
\[
 (-r^2\partial_r^2-3r\partial_r+3)v
 =\eta_Rr^{-3-\delta}+O_\infty(r^{-2-\beta'}).
\]
Variation of constants for the two Euler solutions $r$ and $r^{-3}$,
with the growing solution excluded by decay, yields
\[
 v=r^{-3}v_3(\omega)
   -\frac{\eta_R}{\delta(4+\delta)}r^{-3-\delta}
   +O_\infty(r^{-4}).
\]
Angular difference quotients give a smooth $v_3$, and differentiating
the radial equation gives the stated symbol estimates. In particular
$d_0=\eta_R/[\delta(4+\delta)]>0$.

Define
\[
 \widetilde g=e^{4v}q_R,\qquad \widetilde K=e^{2v}K.
\]
The stress transforms as
$\widetilde K-(\tr_{\widetilde g}\widetilde K)\widetilde g
=e^{2v}(K-(\tr_{q_R}K)q_R)$, so its sign is retained. The conformal
boundary formula is
\[
 H_{\widetilde g}+\tr_T\widetilde K
 =e^{-2v}(H_{q_R}+\tr_TK+4\partial_\nu v),
\]
which preserves the original expansion sign. The scalar and divergence
formulas are
\begin{align}
 e^{4v}\mathcal H_{\widetilde g}(\widetilde K)
 &=\mathcal H_{q_R}(K)+6(e^{4v}-1)
                       -8\Delta_{q_R}v-8|dv|_{q_R}^2,
 \label{eq:str-conformal-hamiltonian}\\
 \mathcal J_{\widetilde g}(\widetilde K)
 &=e^{-2v}\{\mathcal J_{q_R}(K)+4K(\nabla v,\cdot)\}.
 \label{eq:str-conformal-momentum}
\end{align}
The momentum norm in the second formula gains a further $e^{-2v}$.
Substituting Equation~\eqref{eq:str-preparation-equation}, using
$D\sqrt{\eta_R^2+|dv|^2}\geq |K||dv|$, and using the original DEC on
the support of $K$ and its scalar version outside that support, gives
\[
 e^{4v}\{\mathcal H_{\widetilde g}(\widetilde K)
                 -2|\mathcal J_{\widetilde g}(\widetilde K)|\}
 \geq8\eta_Rw_0-8|dv|^2+6(e^{4v}-1)-24v.
\]
The last two terms are nonnegative. Equation~\eqref{eq:str-v-estimate}
and $2\beta>3+\delta$ give
$|dv|^2\leq C\eta_R^2(1+r)^{-3-\delta}$. For large $R$ the required
strict margin follows.

Uniform bilinear convergence of $\widetilde g$ to $g$ gives convergence
of every full-cut infimum: if
$(1-\epsilon)g\leq\widetilde g\leq(1+\epsilon)g$, the same factors
bound all two-dimensional areas, hence their infima. To verify mass
convergence, observe outside $\supp D$ that
\[
 e^{4v}(R_{\widetilde g}+6)
 =\chi_R(R_g+6)+8\eta_Rw_0+O(v^2+|dv|^2).
\]
The weighted integrals of the right side on distant tails are uniformly
small. The scalar-linearization divergence formula in
Lemma~\ref{lem:setup-mass-covariance} has quadratic remainder bounded
by $C V_0r^{-2\beta}$, whose integral is finite because $\beta>3/2$.
Compare the fluxes first on a fixed sphere, where local convergence
applies, and then bound both tails by these integrals. Since
$|V_a|\leq V_0$, this proves convergence for all four fluxes.
Finally, for $K=0$, $D=0$ and the correction equation is linear on
each fixed compact subset for all sufficiently large $R$. Differentiated
local estimates prove the asserted smooth convergence. Relabeling a
sequence $R\to\infty$ proves the lemma.
\end{proof}

\subsection{Full-cut comparison under addition of a short collar}

\begin{lemma}[A collar with almost unchanged cut infimum]
\label{lem:str-added-collar}
Let $(\Omega,g)$ be smooth up to its compact boundary. For every
sufficiently small $\varepsilon>0$ there is an enlarged exterior
$(\Omega_\varepsilon,g_\varepsilon)$, obtained by adding a collar on
the inner side of every boundary component, which agrees with $g$ on
$\Omega$, has a product metric near its new boundary, and satisfies
\begin{equation}
 A_{g_\varepsilon}(\partial\Omega_\varepsilon)
 \geq(1-o(1))A_g(S)\qquad(\varepsilon\downarrow0).
 \label{eq:str-collar-area}
\end{equation}
No constraint inequality on the added collar is required.
\end{lemma}

\begin{proof}
Use disjoint Gaussian collars $[0,L]\times S_j$ with metric
$dx^2+q_j(x)$, where $L>0$ is fixed. Extend $q_j$ smoothly to
$[-\varepsilon,0]$. Since $q_j(x)=q_j(0)+O(\varepsilon)$ there, a
cutoff to the constant $q_j(-\varepsilon)$ near $-\varepsilon$ changes
only the metric values by $O(\varepsilon)$; derivatives may depend on
$\varepsilon$. Choose a smooth map
$F_\varepsilon:[-\varepsilon,L]\to[0,L]$ equal to the identity near
$L$, mapping endpoints to endpoints, and satisfying
$F_\varepsilon-x=O(\varepsilon)$ and
$F_\varepsilon'=1+O(\varepsilon/L)$. Its product with the identity on
$S_j$, extended by the identity elsewhere, is a diffeomorphism
$\Phi_\varepsilon:\Omega_\varepsilon\to\Omega$ with
\[
 (1-o(1))g_\varepsilon\leq\Phi_\varepsilon^*g
                       \leq(1+o(1))g_\varepsilon.
\]
This diffeomorphism maps a connected exterior submanifold and its
entire intrinsic boundary to another of the same type. Applying the
area comparison to every such cut proves
Equation~\eqref{eq:str-collar-area}. The new product face has mean
curvature zero.
\end{proof}

\subsection{The nonlinear Dirichlet problem}
\label{str-pde-section}

For $s\geq0$, define $t(s)\in[0,1)$ by
\begin{equation}\label{str-pde-law}
 s^2=\frac{t}{1-t^{3/2}},\qquad
 b(s)=\frac{s t'(s)}{t(s)}
      =\frac{4(1-t^{3/2})}{2+t^{3/2}}.
\end{equation}
The limiting value of $b$ at zero is $2$. Set
$a(s)=t(s)/s^2$, with $a(0)=1$. Then
\begin{equation}\label{str-pde-law-limits}
\begin{gathered}
 a+s^3a^{3/2}=1,\\
 a=1-s^3+O(s^6),\qquad t=s^2-s^5+O(s^8)
       \qquad(s\downarrow0),\\
 b s^2=\frac{4t}{2+t^{3/2}}\longrightarrow\frac43
       \qquad(s\longrightarrow\infty).
\end{gathered}
\end{equation}
In particular, $0<b\leq2$ and
$\inf_{s\geq0}b(s)(1+s^2)>0$.
The vector $t(|df|_g)\nabla_g f/|df|_g$ is defined to be zero at
critical points. Its continuous dependence on $df$ will always be
understood in this sense.

\begin{theorem}[Fixed-height stress potential]\label{thm:str-stress-dirichlet}
\label{str-pde-fixed-height}
Let $(X,g)$ be a connected smooth complete three-dimensional exterior
with compact smooth nonempty boundary, exactly one end, and a smooth
metric up to its boundary. Let $H_{\partial X}\leq0$, for the normal
pointing into $X$. Suppose that on the end, with $z=r^{-1}$,
\begin{equation}\label{str-pde-prepared-end}
 g=z^{-2}\bar g,\qquad
 \bar g=e^{4v}\left(\frac{dz^2}{1+z^2}+\sigma\right),\qquad
 v=z^3v_3(\omega)-d z^{3+\delta}+O_\infty(z^4),
\end{equation}
where $v_3$ is smooth, $0<\delta<1$, and the remainder and its
tangential derivatives satisfy the indicated symbol estimates.
For every $D>0$ there is a unique bounded weak solution
\begin{equation}\label{str-pde-dirichlet}
 \divg_g\left(t(|df|_g)\frac{\nabla_g f}{|df|_g}\right)=0,
 \qquad f|_{\partial X}=D,\qquad \lim_{r\to\infty}f=0.
\end{equation}
It satisfies $0<f<D$ in the interior, has bounded physical gradient,
belongs to $C^{1,\theta}$ on compact subsets including the finite
boundary for some $\theta>0$, and is smooth wherever $df\ne0$.
On a collar of conformal infinity it is noncritical and
\begin{equation}\label{str-pde-end-expansion}
\begin{gathered}
 f=z C_f(\omega)+R_f,\qquad C_f\in C^\infty(\mathbb S^2),
 \qquad C_f>0,\\
 |(z\partial_z)^j\partial_\omega^\nu R_f|
 \leq C_{j,\nu}z^2\quad(0\leq j\leq2).
\end{gathered}
\end{equation}
The constants can depend on $D$ and the fixed geometry. In particular,
\begin{equation}\label{str-pde-end-flux}
 \begin{gathered}
 |df|_g=zC_f+O(z^2),\qquad
 t(|df|_g)=z^2C_f^2+O(z^3),\\
 -\lim_{r\to\infty}\int_{\{r=\mathrm{const}\}}
 t(|df|_g)\left\langle\frac{\nabla_g f}{|df|_g},\nu_g\right\rangle
 dA_g=\int_{\mathbb S^2}C_f^2\,dA_\sigma .
 \end{gathered}
\end{equation}
\end{theorem}

We give the a priori estimates and the existence argument separately.
In what follows, constants are independent of the outer truncation
and of the regularization parameter, unless stated otherwise.

\begin{lemma}[Regularization and the public regularity input]
\label{str-pde-regularity-input}
For sufficiently small $h>0$, choose a smooth $b_h$ equal to $1$ on
$[0,h/2]$, equal to $b$ on $[2h,\infty)$, and between $1$ and $2$ on
the intervening interval. Define $t_h$ by
$s t_h'/t_h=b_h$ and $t_h(2h)=t(2h)$. For each fixed $M<\infty$,
\begin{equation}\label{str-pde-regularized-structure}
 c_M(h+s)s\leq t_h(s)\leq C_M(h+s)s,\qquad
 c_M(h+s)\leq t_h'(s)\leq C_M(h+s)
 \quad(0\leq s\leq M).
\end{equation}
Here $t_h(s)/s$ is constant near zero, so the regularized vector
field is smooth at the zero covector. On coordinate patches with
uniformly controlled metric and first derivatives, its density
\[
 A_h^i(x,p)=\sqrt{\det g}\,\frac{t_h(|p|_g)}{|p|_g}g^{ij}p_j
\]
obeys, for $|p|\leq M$,
\begin{equation}\label{str-pde-density-structure}
 \begin{split}
 \lambda_M(h^2+|p|^2)^{1/2}|\zeta|^2
 &\leq D_pA_h[\zeta,\zeta],\qquad
 |D_pA_h|\leq\Lambda_M(h^2+|p|^2)^{1/2},\\
 |A_h|+|\partial_x A_h|&\leq C_M(h+|p|)|p|.
 \end{split}
\end{equation}
Once an a priori gradient bound $M$ has been established, the field
can be extended for $|p|>M$ to satisfy these inequalities globally,
with possibly changed constants independent of $h$.

Consequently bounded weak solutions with this gradient bound have
local $C^{1,\theta}$ estimates, including Dirichlet boundary patches,
uniformly in $h$. On interior unit patches for $h=0$, the estimate
can be normalized to give
\begin{equation}\label{str-pde-oscillation-gradient}
 \sup_{B_{1/2}}|df|_g\leq C\operatorname{osc}_{B_1} f,
\end{equation}
provided the same fixed extension and the same structural constants
apply on the patch. Smooth boundary data on finite boundary patches
are allowed.
\end{lemma}
\begin{proof}
Below $h/2$ the normalization gives $t_h(s)=c_hs$ with
$c_h\asymp h$. On the transition interval $s\asymp h$,
$t_h\asymp h^2$, and $b_h\in[1,2]$. Above it the inequalities
follow from $t/s^2>0$ and $b>0$ on a bounded interval.
The radial and tangential eigenvalues of the covector derivative
are $t_h'$ and $t_h/s$. Metric differentiation costs at most a
constant times $|p|_g$ in differentiating $|p|_g$, and
$s t_h'=b_h t_h\leq2t_h$. These facts prove
\eqref{str-pde-density-structure}.
For example, outside a slightly larger bounded-gradient interval
one may interpolate the logarithmic derivative to $2$ and then
continue by a quadratic flux law. Positive lower and upper bounds
for that logarithmic derivative give the required global extension.
This extension is used only after the gradient bound; it is not an
existence argument for the saturating law.

We specify the hypotheses of the regularity theorem.
For $-\partial_i A^i(x,u,Du)=B(x,u,Du)$ on a $C^{1,\alpha}$
domain, Appendix Theorem A.1 of \cite{porrettaveron2009} permits
\[
 \lambda(\mu^2+|p|^2)^{(p_0-2)/2}I
 \leq D_pA\leq
 \Lambda(\mu^2+|p|^2)^{(p_0-2)/2}I,
\]
an analogous bound for $|D_pA|$, and the spatial and scalar
continuity condition
\[
 |A(x,s,p)-A(y,t,p)|
 \leq C(1+|p|^{p_0-2}+|p|^{p_0-1})
       (|x-y|^\alpha+|s-t|^\alpha),
 \qquad |B|\leq C(1+|p|^{p_0}).
\]
For bounded weak solutions it yields a H\"older continuous gradient
up to a $C^{1,\alpha}$ Dirichlet boundary; the same local argument
gives the interior estimate. We apply it with $p_0=3$, $\mu=h$,
no scalar dependence, and $B=0$. The metric bounds and
\eqref{str-pde-density-structure} verify every displayed condition;
constants on a fixed height and gradient range are uniform in $h$.
At a finite Dirichlet face the datum is constant; subtracting that
constant reduces a boundary patch to the homogeneous theorem.
The finitely many smooth boundary components are disjoint, so
these patches can be chosen separately. At conformal infinity
the datum is already zero.
The interior gradient estimates originate in
\cite{tolksdorf1984}, and the frozen boundary estimate used in the
proof is the boundary regularity estimate of \cite{lieberman1988}.
We use the stated Appendix theorem as the precise boundary input.

For completeness, normalization removes any additive constant in
the interior gradient bound. Extend the unregularized density with
$D_pA\asymp|p|$ and $|\partial_xA|\leq C|p|^2$. If
$m=\operatorname{osc}_{B_1}f>0$, subtract the infimum and put
$u=f/m$. The equation for $u$ has density
$A_m(x,p)=m^{-2}A(x,mp)$, whose structural constants do not depend
on $m$. Since $0\leq u\leq1$, the local gradient estimate for $u$
is uniform. Multiplication by $m$ proves
\eqref{str-pde-oscillation-gradient}. The case $m=0$ is immediate.
\end{proof}

\begin{lemma}[Boundary barriers and global gradient bound]
\label{str-pde-global-gradient}
Let $f$ be a classical regularized solution on a truncation $X_R$,
with boundary values $D'$ at the finite boundary and zero at
$r=R$, where $0\leq D'\leq D$. Then
\begin{equation}\label{str-pde-global-bounds}
 0\leq f\leq D,\qquad |df|_g\leq M_D,
 \qquad f\leq C_Dr^{-\xi}\quad(r\geq r_1),
\end{equation}
for any fixed $0<\xi<1$ and suitable $r_1,C_D,M_D$ independent
of small $h$ and sufficiently large $R$.
\end{lemma}
\begin{proof}
At a noncritical point put $s=|df|_g$, $e=\nabla f/s$, and
\[
 Lf=A^{ij}\nabla_i\nabla_jf=0,
 \qquad A=g^{-1}+(b_h-1)e\otimes e.
\]
The maximum principle gives the height bound.
Let $d$ be the inward distance from the finite boundary. Its level
sets have mean curvature $H(d)=\Delta_g d\leq C d$ in a fixed
collar, since $H(0)\leq0$. Consider
\[
 \underline f=D'-\gamma(d),\qquad
 \gamma'(d)=\frac{c}{\sqrt{d^2+d_*^2}}.
\]
The radial operator on this function is
\begin{equation}\label{str-pde-inner-barrier}
 L\underline f
 =\gamma'\left(\frac{b_h(\gamma')d}{d^2+d_*^2}-H(d)\right).
\end{equation}
Uniformly in $h$, $b_h(s)(1+s^2)\geq c_0>0$, and hence
\[
 \frac{b_h(\gamma')}{d^2+d_*^2}
 \geq\frac{c_0}{d^2+d_*^2+c^2}.
\]
Choose the collar length $\ell$ and $c>0$ small first; then choose
$d_*>0$ sufficiently small. The last bound makes
\eqref{str-pde-inner-barrier} nonnegative on $[0,\ell]$, while
$\gamma(\ell)=c\operatorname{arsinh}(\ell/d_*)\geq D$.
Comparison with $\underline f$ and the constant upper barrier
gives a uniform finite-boundary gradient bound $c/d_*$.

On the end, for $F(r)=r^{-\xi}$, the leading radial operator is
$\xi(b_h\xi-2)r^{-\xi}$, with relative error tending to zero
uniformly for $0<b_h\leq2$. Thus $F$ is a strict supersolution
beyond a fixed sphere, uniformly in any positive multiple of $F$.
Comparison with
$K_D(r^{-\xi}-R^{-\xi})$ gives both the last estimate in
\eqref{str-pde-global-bounds} and the outer-boundary gradient
bound. Choose $K_D$ at the fixed inner sphere before sending
$R$ to infinity; for $R$ at least twice its radius the required
$K_D$ is uniform.

It remains to exclude a large interior gradient. The following
calculation is for $s\geq s_0$, where $b_h=b$.
Writing $x=t^{3/2}$ gives
\[
 b+s b'=b\left(1-\frac{18x}{(2+x)^2}\right).
\]
Thus, for fixed sufficiently large $s_0$, $b<1$ and $b+s b'<0$.
At a point with first frame vector $e$, differentiating $Lf=0$
and commuting covariant derivatives yields the exact identity
\begin{equation}\label{str-pde-log-gradient}
\begin{aligned}
 L\log s
 &=\frac{1}{s^2}\Bigl(
 |\Hess f|_{e^\perp\times e^\perp}^2
 +(1-b)|\Hess f(e,\cdot)|_{e^\perp}^2\Bigr)\\
 &\quad-\frac{b+s b'}{s^2}\Hess f(e,e)^2+\Ric_g(e,e).
\end{aligned}
\end{equation}
To see the coefficients, differentiate $A$ in the $e$ direction:
its contraction with $\Hess f$ is
$b'\Hess f(e,e)^2+
 2(b-1)s^{-1}|\Hess f(e,\cdot)|_{e^\perp}^2$.
The two derivatives of $\log s$ then leave exactly the three
quadratic terms displayed. In the curvature commutator the
additional $(b-1)e\otimes e$ term vanishes by antisymmetry, leaving
$\Ric(e,e)$.
The Ricci curvature is bounded below on the whole exterior, so
$L\log s\geq-C$. Moreover $L(f^2)=2b s^2$ and
$\inf_{s\geq s_0}b s^2>0$. Choose $c_1$ so large that
\[
 L(\log s+c_1f^2)>0\qquad(s\geq s_0).
\]
An interior maximum in this range is impossible. The two boundary
gradient bounds and $0\leq f\leq D$ therefore bound the maximum
of $\log s+c_1f^2$, and then $s$, everywhere. At a critical point
$\log s=-\infty$, so such points cause no difficulty in this
maximum argument.
\end{proof}

\begin{proof}[Existence, uniqueness, and finite-boundary regularity in
Theorem~\ref{str-pde-fixed-height}]
Fix $h>0$ and a finite smooth truncation. Use the homotopy of
boundary values $\lambda D$, $0\leq\lambda\leq1$, for the original
regularized saturating equation. At $\lambda=0$ the zero solution
exists. For fixed $h$ and on the gradient range of
Lemma~\ref{str-pde-global-gradient}, the operator is uniformly
elliptic and smooth in the covector. Its linearization has no
zeroth-order term, is invertible with homogeneous Dirichlet data
by the maximum principle and linear elliptic solvability, and
gives openness. For closedness, first apply
Lemma~\ref{str-pde-regularity-input} using the already established
gradient bound. Its extension agrees with the original flux on
every solution. This gives a common $C^{1,\theta}$ bound; the
ordinary uniformly elliptic Schauder estimates for this fixed
$h$ then give $C^{2,\theta'}$ bounds and the required compactness.
The continuity method therefore reaches $\lambda=1$; the linear
Dirichlet solvability and Schauder estimates used here follow from
Lemma~\ref{lem:linear-separated-faces}, including its general-principal
all-Dirichlet conclusion.
No uniform ellipticity as $h\downarrow0$ is asserted at this step.

Now use the estimates that are uniform in $h$ and in the outer
truncation. On every fixed interior or finite-boundary patch,
the $C^{1,\theta}$ estimates permit a diagonal $C^1$ limit as
$h\downarrow0$ and $R\to\infty$. The fluxes converge locally
uniformly, so the weak equation and finite Dirichlet trace pass
to the limit. The common bound $C_Dr^{-\xi}$ gives the zero
limiting value at infinity. Interior elliptic bootstrap applies
on every set where $df\ne0$.

The unregularized flux is strictly monotone in the covector:
its radial potential has derivative $t(s)$, with $t'>0$ for
$s>0$. For two solutions $f_1,f_2$, test the difference of the
weak equations by $(f_1-f_2-\varepsilon)_+$. This test has compact
support, since both functions tend uniformly to zero at infinity,
and has zero trace on the finite boundary. Strict monotonicity
forces its gradient to vanish. It is therefore zero; letting
$\varepsilon\downarrow0$ proves $f_1\leq f_2$, and interchanging
the solutions proves uniqueness. The same comparison proves
monotonicity in $D$.
Positivity can be propagated by a direct local comparison.
On a compact geodesic annulus $r_0<\varrho<r_1$ lying below
the injectivity radius, take
$w=A(e^{-k\varrho^2}-e^{-kr_1^2})$. Its radial operator is
$b w''+w'\Delta_g\varrho$. First choose $k$ so large that
$b_0(2k\varrho-1/\varrho)\geq\Delta_g\varrho$ on the annulus,
where $b_0=\min_{0\leq s\leq1}b(s)>0$. Then take $A>0$
small enough that $|dw|\leq1$ and the inner boundary value
of $w$ is below a given positive minimum of $f$. Since
$w''/(-w')=2k\varrho-1/\varrho$, this is a positive
subsolution, zero at the outer sphere. Weak comparison
propagates positivity from the inner ball to the larger
ball. The finite-boundary barrier supplies an initial
positive collar; a finite chain of interior balls along
any path proves $f>0$ throughout the connected interior.
Apply the same argument to $D-f$, which solves the same
odd-flux equation and is positive near infinity, to obtain
$f<D$. This completes the asserted existence and
finite-boundary statements.
\end{proof}

\subsection{The conformal boundary and its expansion}

\begin{lemma}[Linear height and physical-gradient bounds at infinity]
\label{str-pde-linear-end-bounds}
For the solution in Theorem~\ref{str-pde-fixed-height} there are
$c,C,Q>0$ and a collar $0<z<z_0$ such that
\begin{equation}\label{str-pde-linear-bounds}
 cz\leq f\leq Cz,\qquad |df|_g\leq Qz.
\end{equation}
These constants can be chosen uniformly for a family whenever
the coarse bounds $f\leq C_0z^\xi$, a positive lower bound on one
fixed inner sphere, local physical gradient bounds, and the
compact metric bounds used in the proof are uniform.
\end{lemma}
\begin{proof}
There is a useful radial identity that also applies to the
families considered below. For $g=z^{-2}\bar g$, put
$u=|dz|_{\bar g}^2$ and
$\bar e=\nabla_{\bar g}z/\sqrt u$. If $F'>0$, the operator on
$F(z)$, divided by $z^2u$, is
\begin{equation}\label{str-pde-compact-radial}
 bF''+\left(\frac{b-2}{z}+E\right)F',\qquad
 E=\frac{\Delta_{\bar g}z+(b-1)\Hess_{\bar g}z(\bar e,\bar e)}u,
 \qquad b=b(z\sqrt u F').
\end{equation}
Indeed
$\Hess_g z=\Hess_{\bar g}z+
 2z^{-1}dz\otimes dz-z^{-1}u\bar g$ and
$\Delta_g z=z^2\Delta_{\bar g}z-zu$.
For \eqref{str-pde-prepared-end},
\[
 E=\frac{b z}{1+z^2}-2(b-2)v_z=O(z)
\]
uniformly for $0<b\leq2$. More generally, if
$\bar g-h_0=O_2(z^\beta)$ with
$h_0=(1+z^2)^{-1}dz^2+\sigma$ and $\beta>1$, then
$E=O(z)+O(z^{\beta-1})$, again uniformly in $b$.

Choose $F_+(z)=A(z-c_2z^2)$ with a fixed $c_2>0$.
The term $(b-2)F_+'/z$ is nonpositive. Make $z_0$ small
and take $A$ sufficiently large for $F_+(z_0)\geq f(z_0,\cdot)$.
The order of these choices can be made consistent with small
physical slope: the coarse bound allows
$A=2C_0z_0^{\xi-1}$, so $z\sqrt u F_+'\leq C C_0z_0^\xi$.
Thus $b\geq1$ on the whole collar after reducing $z_0$.
The negative term $-2Ac_2b$ then dominates $E F_+'$.
Hence $F_+$ is a supersolution, and comparison gives $f\leq F_+$.
For the lower barrier take $F_-(z)=a_0(z+c_2z^2)$, with $a_0>0$
small enough to lie below the positive minimum of $f$ on
$z=z_0$. Here $b=2-O(a_0^3z^3)$, so the unfavorable term
$(b-2)F_-'/z$ is $O(a_0^4z^2)$, while $bF_-''$
has a positive constant multiple of $a_0$. Reducing the collar
if needed makes this a subsolution. For the upper comparison test by $(f-F_+-\varepsilon)_+$;
for the lower comparison use $(F_--f-\varepsilon)_+$.
The limiting zero boundary values make both tests compactly
supported for every $\varepsilon>0$, and their traces vanish
on the inner sphere. Monotonicity and then
$\varepsilon\downarrow0$ justify both comparisons at infinity.
They prove the height bounds in \eqref{str-pde-linear-bounds}.

Physical balls of a fixed small radius in this collar have
uniform metric bounds; $z$ is comparable to its central value
throughout each ball. The height estimate gives oscillation
$O(z)$ there. Apply \eqref{str-pde-oscillation-gradient} using
the bounded-gradient extension furnished by the global bound
(or the asserted uniform local bound for a family). It gives
$|df|_g=O(z)$. All the choices above use only the quantities
listed in the statement, proving the uniform version.
\end{proof}

\begin{lemma}[Compactified equation and boundary jet]
\label{str-pde-boundary-jet}
Under \eqref{str-pde-prepared-end} and
\eqref{str-pde-linear-bounds}, the expansion
\eqref{str-pde-end-expansion} holds.
\end{lemma}
\begin{proof}
Put $q=|df|_{\bar g}$. In dimension three, change of the volume
density and inverse metric gives the exact equation
\begin{equation}\label{str-pde-compact-equation}
 \divg_{\bar g}\big(a(zq)q\nabla_{\bar g}f\big)=0.
\end{equation}
For reference, the conformal factor in
\eqref{str-pde-prepared-end} can also be cancelled exactly:
with $Q=|df|_{h_0}$ the equation is
$\divg_{h_0}[a(ze^{-2v}Q)Q\nabla_{h_0}f]=0$.
The bounds \eqref{str-pde-linear-bounds} give $q\leq Q_0$,
so $f$ is Lipschitz up to $z=0$ in the compact metric and has
zero trace there. In a $\bar g$-orthonormal frame, the covector
derivative of the compact flux has eigenvalues
\[
 a(zq)q,\qquad a(zq)q\,b(zq).
\]
They are uniformly comparable to $q$, and the first spatial
derivatives of the coordinate-density flux are bounded by
$Cq^2$ on this bounded covector range. A continuation beyond
that range gives the global three-Laplace structure used in
Lemma~\ref{str-pde-regularity-input}.

The interior equation has the Dirichlet weak formulation up
to $z=0$. Indeed, cut a smooth zero-trace test off in
$z<\varepsilon$. The compact flux is bounded, the test is
$O(z)$, and the cutoff layer has volume $O(\varepsilon)$;
the additional integral is therefore $O(\varepsilon)$.
No assertion of zero normal flux is made.
The boundary gradient estimate now yields $f\in C^{1,\theta}$
up to $z=0$. The lower bound $f\geq cz$ implies
$f_z(0,\omega)\geq c$. Thus a smaller full collar has gradient
bounded away from zero and the compact equation is uniformly
elliptic there.

The compact metric is generally $C^{3,\delta}$, rather than
smooth in its normal variable. Every one of its angular
derivatives has the same finite normal regularity. Choose
$0<\gamma<\min(\theta,\delta)$. The first boundary upgrade can be
obtained without assuming a second derivative in advance. Take angular
difference quotients of the divergence equation. Their principal
coefficient is the average of $A_p$ between two gradients; it has a
common H\"older modulus and is uniformly elliptic. Their divergence
source is the difference quotient of the explicit spatial dependence
of the flux, with the same H\"older control. The boundary trace remains
zero. The linear divergence boundary estimate
\cite[Theorem~9.1, p.~675, with $l=1$]{agmondouglisnirenberg1959}
therefore bounds these difference quotients in $C^{1,\gamma}$ on smaller
half-balls, using their already bounded $C^0$ norm. Passing to the limit
gives tangential second derivatives and mixed derivatives.
The equation then gives the normal derivative of the normal flux
from its tangential derivatives. Its derivative with respect to $f_z$
is bounded away from zero, so the implicit function theorem recovers
the remaining normal second derivative with its H\"older bound.
Thus $f\in C^{2,\gamma}$, and subsequent ordinary Schauder estimates
use Lemma~\ref{lem:linear-separated-faces}.
To record explicitly why arbitrary angular derivatives are
allowed, write \eqref{str-pde-compact-equation} in density form
$\partial_j A^j(x,Df)=0$ and put
$B^{jk}=A^j_{p_k}(x,Df)$. For an angular multi-index $\nu\ne0$,
$w_\nu=\partial_\omega^\nu f$ satisfies
\begin{equation}\label{str-pde-tangential-equation}
 \partial_j(B^{jk}\partial_k w_\nu)=-\partial_j F_\nu^j,
 \qquad w_\nu|_{z=0}=0.
\end{equation}
After the top derivative $B Dw_\nu$ is removed, the chain-rule
remainder $F_\nu$ contains only explicit angular derivatives
of the flux and $D\partial_\omega^\mu f$ with $|\mu|<|\nu|$.
Inductively $B,F_\nu\in C^{1,\gamma}$, and the preceding
stage gives $w_\nu\in C^{1,\gamma}$. Expansion of
\eqref{str-pde-tangential-equation} gives a uniformly elliptic
nondivergence equation whose right side is $C^{0,\gamma}$.
Linear boundary Schauder estimates give
$w_\nu\in C^{2,\gamma}$. Finite covers and nested boundary
patches complete this induction for each fixed $\nu$, with
constants allowed to depend on $\nu$.

Set $C_f=f_z|_{z=0}$. It is positive and smooth angularly.
Taylor's formula applied to each $w_\nu$ gives
$\partial_\omega^\nu(f-zC_f)=O(z^2)$. The identities
\[
 z\partial_z(w_\nu-z(w_\nu)_z|_0)
 =z((w_\nu)_z-(w_\nu)_z|_0),
\]
\[
 (z\partial_z)^2(w_\nu-z(w_\nu)_z|_0)
 =z((w_\nu)_z-(w_\nu)_z|_0)+z^2(w_\nu)_{zz}
\]
prove the two symbol derivative estimates. Finally
$q=C_f+O(z)$, and the outward normal to the physical end points
toward decreasing $z$. Substituting in the physical flux, with
$dA_g=z^{-2}e^{4v}dA_\sigma$, gives
\eqref{str-pde-end-flux}. This proves the remaining assertions
of Theorem~\ref{str-pde-fixed-height}.
\end{proof}

\subsection{Uniform estimates when the finite-boundary heights grow}

The preceding global gradient bound depends on $D$. The following
local estimate supplies the different uniformity needed when
$D$ tends to infinity.

\begin{lemma}[A gradient estimate from local height]
\label{str-pde-local-gradient}
On a physical coordinate ball with uniformly bounded geometry,
a solution with $0\leq f\leq H$ has a gradient bound on a
concentric smaller ball depending only on $H$ and that geometry.
The estimate applies uniformly to smooth solutions off their
critical sets and to their $C^1$ limits.
\end{lemma}
\begin{proof}
Increase $s_0$ in \eqref{str-pde-log-gradient} if necessary.
Since $b\to0$ and $-(b+s b')/b\to1$, that identity implies
\begin{equation}\label{str-pde-strong-log-gradient}
 L\log s\geq\tfrac12|d\log s|_A^2-C
 \qquad(s\geq s_0).
\end{equation}
Choose a smooth positional cutoff $\psi$ that is negative near
the ball boundary and positive on the smaller ball, with
controlled first two derivatives. Choose $\varepsilon>0$
depending only on $H$ so that
$\phi=\psi+\varepsilon f$ is still negative near the outer
boundary and uniformly positive on the smaller ball.
Because $Lf=0$, $L\phi=L\psi$ is uniformly bounded. Moreover,
when $s$ is sufficiently large,
\begin{equation}\label{str-pde-axial-cutoff}
 |d\phi|_A^2\geq b(\varepsilon s+e\psi)^2
 \geq c\varepsilon^2,
\end{equation}
using $bs^2\to4/3$.
On the relatively compact region $\phi>0$, put
$q(\phi)=e^{k\phi}-1$ and consider
$W=\log s+4\log q(\phi)$.
At an interior maximum with large gradient,
$d\log s=-4(q'/q)d\phi$. Equations
\eqref{str-pde-strong-log-gradient}--\eqref{str-pde-axial-cutoff}
then give
\begin{equation}\label{str-pde-local-maximum}
 LW\geq
 4\left((q'/q)^2+q''/q\right)|d\phi|_A^2
 +4(q'/q)L\phi-C.
\end{equation}
Here $q'/q\geq k$ and $q''/q=kq'/q$. A fixed sufficiently
large $k$ makes the right side positive, a contradiction.
Since $W\to-\infty$ at $\phi=0$, and $q$ is bounded below
on the smaller ball, a bound at the remaining, bounded-gradient
maxima gives the desired interior gradient bound. Every
calculation takes place where the gradient is large and the
solution is smooth, so critical points introduce no additional
regularity requirement.
\end{proof}

\begin{lemma}[Common end for growing heights]
\label{lem:str-uniform-end}
\label{str-pde-growing-heights}
Let $g_i$ be smooth metrics on the same exterior, with
$H_{\partial X}(g_i)\leq0$, converging smoothly on compact
subsets to a smooth metric $g$. Suppose that each has the
prepared end required for Theorem~\ref{str-pde-fixed-height},
and that for a common $\beta>1$ and $M<\infty$,
\begin{equation}\label{str-pde-uniform-physical-metrics}
 \sum_{j=0}^2|\nabla_b^j(g_i-b)|_b\leq Mz^\beta
\end{equation}
on a common end. Assume the compact metrics are uniformly
positive definite. Let $D_i\geq1$ be arbitrary finite heights
and let $f_i$ be their solutions. Then on a common end collar
there are constants $c,C,Q>0$, independent of $i,D_i$, such that
\begin{equation}\label{str-pde-uniform-linear-bounds}
 cz\leq f_i\leq Cz,\qquad |df_i|_{g_i}\leq Qz.
\end{equation}
For some $\gamma>0$ the functions have uniform compact
$C^{2,\gamma}$ bounds on a common smaller collar, which is
noncritical for every $f_i$. After extraction, they converge
in $C^{2,\gamma'}$ for $0<\gamma'<\gamma$ on a still smaller
closed compact collar, and smoothly on interior compact subsets
of that collar, to a solution for $g$.
If $g$ has a smooth compactification in another defining
function $\rho$ comparable to $z$, the limiting solution is
smooth there and has $f=\rho C_\infty+O(\rho^2)$ with positive
smooth $C_\infty$.
\end{lemma}
\begin{proof}
First obtain an upper height bound without using $D_i$.
Write $\eta=\operatorname{arsinh}r$ and choose a large fixed
$\eta_0$. On $x=\eta-\eta_0>0$ take
\begin{equation}\label{str-pde-pole-barrier}
 w(\eta)=\frac{A e^{-\xi x}}x,\qquad 0<\xi<1.
\end{equation}
Uniformly in $i$, $|d\eta|_{g_i}^2=1+o(1)$ and
\[
 \Delta_{g_i}\eta+(b-1)\Hess_{g_i}\eta(e,e)=2+o(1),
 \qquad e=\frac{\nabla_{g_i}\eta}{|d\eta|_{g_i}},
\]
for $0<b\leq2$. The radial operator is therefore
$b w''|d\eta|^2+w'(2+o(1))$. A direct differentiation gives
\[
 \frac{w''}{-w'}=\xi+\frac1x+\frac1{x(1+\xi x)}.
\]
Far from the pole, $b\leq2$ and $\xi<1$ give a strict
supersolution after fixing $\eta_0$ sufficiently large.
Near the pole, $|w'|\asymp A/x^2$, $w''\asymp A/x^3$,
and $b(|w'||d\eta|)=O(x^4/A^2)$: the positive term is
$O(x/A)$, whereas the negative term has order $-A/x^2$.
On the remaining compact interval, increasing $A$ makes $b$
uniformly small. Thus one fixed $A$ makes $w$ a supersolution
on the entire end, for every $i$.

For each finite $D_i$ compare first on a truncated domain
starting sufficiently near the pole that $w>D_i$. At its
outer boundary the Dirichlet solution is zero. Passage to
the exhaustion limit gives $f_i\leq w$. In particular
$f_i\leq C_0z^\xi$ beyond a fixed sphere, independently of
$i,D_i$. Applying Lemma~\ref{str-pde-local-gradient} on
uniform physical balls there gives a common local physical
gradient bound. This is the bound needed before extending
the flux to a three-Laplace structure law on these patches.

For a uniform positive lower value on a fixed sphere, let
$u_i$ be the height-one solution for $g_i$. Comparison gives
$f_i\geq u_i$. The fixed-height estimates have uniform
constants in this family: local smooth convergence controls
the finite-boundary collars and compact geometry, while
\eqref{str-pde-uniform-physical-metrics} controls the end.
Thus every subsequence has a further subsequence converging
locally in $C^1$ to the height-one solution for $g$, with
zero limiting boundary value at infinity. This limit is
positive in the interior by the local annular comparison
used above. Uniqueness gives the same limit for every
subsequence, so the minima of $u_i$ on a fixed interior
sphere have a positive common lower bound. Reducing it to
handle finitely many initial terms if necessary supplies
the desired lower bound.
Lemma~\ref{str-pde-linear-end-bounds} now gives
\eqref{str-pde-uniform-linear-bounds}, using its general
radial formula \eqref{str-pde-compact-radial}.

We spell out the regularity conversion for the metrics.
If $E_i=z^2(g_i-b)$, tensor-to-coordinate conversion in
\eqref{str-pde-uniform-physical-metrics} gives
\begin{equation}\label{str-pde-weight-conversion}
 |\partial^j E_i|\leq C M z^{\beta-j},\qquad 0\leq j\leq2.
\end{equation}
Indeed, a tensor of rank $2+j$ with $b$-norm $O(z^\beta)$
has coordinate components $O(z^{\beta-2-j})$, and
$\Gamma_b=O(z^{-1})$, $\partial\Gamma_b=O(z^{-2})$.
The compact metrics are consequently uniformly $C^{1,\alpha}$
for any $0<\alpha<\min(\beta-1,1)$. To verify the only
nontrivial case $1<\beta<2$, let two points be separated by
coordinate distance $d$. If $d$ is smaller than half the
larger $z$ value, integrate the second derivative bound
along a path on which $z$ is comparable to that value:
\[
 |DE_i(x)-DE_i(y)|\leq C d z^{\beta-2}
 \leq C d^{\beta-1}.
\]
For the remaining pairs use the first derivative amplitude
$Cz^{\beta-1}$ at the two endpoints. Pairs on the boundary
follow by continuity. Smooth interior convergence and this
uniform estimate also give compact $C^{1,\alpha'}$
convergence for $\alpha'<\alpha$. In particular, a cutoff
annulus moving toward $z=0$ does not require any extra
ordinary-derivative bound.

The compact flux in \eqref{str-pde-compact-equation} has
uniform three-Laplace structural constants, and
\eqref{str-pde-uniform-linear-bounds} gives a common compact
gradient bound. The weak zero-trace formulation and
Lemma~\ref{str-pde-regularity-input} give uniform
$C^{1,\theta}$ estimates up to $z=0$. The lower height
bound implies $f_{i,z}(0,\omega)\geq c$. The H\"older
gradient estimate therefore makes a common collar
noncritical. On it ordinary quasilinear boundary Schauder
estimates yield a uniform $C^{2,\gamma}$ bound for
$0<\gamma<\min(\alpha,\theta)$. Compactness gives the
stated convergence, and smooth interior convergence of
the metrics permits elliptic bootstrap on each compact
subset of the collar away from $z=0$.

Finally, in a smooth limiting compactification
$\widehat g=\rho^2g$, comparability of $\rho$ and $z$ gives
$c'\rho\leq f\leq C'\rho$ and
$|df|_{\widehat g}\leq Q'$. Apply the compactified boundary
argument afresh with $\rho$. It gives a noncritical collar;
all subsequent boundary Schauder bootstraps are now allowed
because the limiting compact metric is smooth. Hence the
limit is smooth up to that conformal boundary.
\end{proof}

\begin{remark}\label{str-pde-uniformity-scope}
The uniform physical control through two derivatives in
Lemma~\ref{str-pde-growing-heights} gives finite compact
regularity. It does not give uniform bounds for arbitrary
angular derivatives of the individual coefficients $C_{f_i}$.
Each fixed prepared metric has the angular smoothness in
\eqref{str-pde-end-expansion}, with constants depending on
that metric. Smoothness of the limit in the final assertion
is obtained separately from the smooth limiting
compactification. The later application has $\beta>3/2$,
which is stronger than the $\beta>1$ needed for this
regularity conversion.
\end{remark}

\subsection{Flux saturation and the conserved stress}

\begin{lemma}[Flux saturation for full enclosing cuts]
\label{lem:str-flux-saturation}
Let $f_D$ be the solution in Theorem~\ref{thm:str-stress-dirichlet}
on an exterior $(X,g)$ with inner boundary $T$, with
$f_D|_T=D>0$ and $f_D\to0$ at infinity. Write
$f_D=zC_D(\omega)+O(z^2)$, and set
\[
 \mathcal F_D=\int_{\mathbb S^2} C_D^2\,d\omega.
\]
Then
\begin{equation}
 \int_X t(|df_D|)|df_D|\,dV_g=D\mathcal F_D,
 \qquad 0\leq\mathcal F_D\leq A_g(T),
 \qquad \lim_{D\to\infty}\mathcal F_D=A_g(T).
 \label{eq:str-flux-saturation}
\end{equation}
In particular
\begin{equation}
 \fint_{\mathbb S^2}C_D^3\,d\omega
 \geq\left(\frac{\mathcal F_D}{4\pi}\right)^{3/2}.
 \label{eq:str-jensen}
\end{equation}
\end{lemma}

\begin{proof}
Put $X_D=t(|df_D|)\nabla f_D/|df_D|$, with value zero at a critical
point. This vector field is continuous, divergence free in
distributions, and has norm at most one. The end expansion gives
$|df_D|=zC_D+O(z^2)$ and $t(|df_D|)=z^2C_D^2+O(z^3)$. Its flux
through a distant sphere, with outward normal, tends to
$-\mathcal F_D$. The divergence theorem therefore gives the inward
flux $\mathcal F_D$ on every full cut separating the inner boundary
from infinity. Consequently $\mathcal F_D$ is bounded by the area of
every admissible cut. This is also valid for continuous divergence-free
fields: apply the distributional equation to smooth approximations of
the indicator of the region between a cut and a distant sphere, or
first mollify the density on that compact region.

Integration of $\operatorname{div}(f_DX_D)=t(|df_D|)|df_D|$ over a
truncation gives the first identity in
Equation~\eqref{eq:str-flux-saturation}. The outer boundary term tends
to zero because $f_D=O(z)$ while the flux has a finite limit. The
integrand is integrable at infinity: it is $O(z^3)$ and the volume
element is $O(z^{-3})\,dz\,d\omega$.

Fix $0<\lambda<1$. The coarse radial supersolution in
Theorem~\ref{thm:str-stress-dirichlet} gives
$f_D\leq cD r^{-\xi}$ for a fixed $0<\xi<1$, with $c$ independent
of $D$. Hence
\begin{equation}
 \Vol_g\{f_D>\lambda D\}\leq V_\lambda<\infty
 \label{eq:str-superlevel-volume}
\end{equation}
uniformly in $D$. The function $\sigma(1-t(\sigma))$ is bounded on
$[0,\infty)$: it tends to zero both at zero and at infinity, since
$1-t(\sigma)=O(\sigma^{-2})$ at infinity. Thus for a constant $C_0$,
\[
 D\mathcal F_D
 \geq\int_{\{f_D>\lambda D\}}|df_D|\,dV_g-C_0V_\lambda.
\]

We justify the coarea lower bound with the exact full-boundary
convention. For almost every level $a\in(0,D)$, the set
$\{f_D>a\}$ has finite perimeter, contains an inner collar of $T$,
and is bounded. Its complement has a unique component containing the
end. Fill the other complementary components. This operation can only
decrease the perimeter, and the remaining entire boundary separates
the end from every component of $T$. Approximate the filled set by
smooth sets in a fixed compact region while preserving a smaller
inner collar and the distant exterior; their perimeter converges to
the finite-perimeter value. Such approximation follows by convolving
the characteristic function in finitely many interior charts and
using regular levels; the inner and outer collars are held fixed, and
bounded complementary components are again filled. Their full
boundaries are smooth admissible cuts. It follows that the original
level perimeter is at least $A_g(T)$. Coarea, valid for the locally
Lipschitz $f_D$, now gives
\[
 \int_{\{f_D>\lambda D\}}|df_D|\,dV_g
 =\int_{\lambda D}^{D}\operatorname{Per}\{f_D>a\}\,da
 \geq(1-\lambda)D A_g(T).
\]
No smooth Sard assertion at critical points of a $C^{1,\theta}$
solution is needed. We obtain
\[
 (1-\lambda)A_g(T)-C_0V_\lambda/D
 \leq\mathcal F_D\leq A_g(T).
\]
Let $D\to\infty$, then $\lambda\downarrow0$. Finally,
Equation~\eqref{eq:str-jensen} is H\"older's Inequality on the sphere
with its probability measure $d\omega/(4\pi)$.
\end{proof}

\begin{lemma}[The conserved stress and its smoothing]
\label{lem:str-stress-identities}
For a solution $f$ of the saturating equation, put
$\sigma=|df|$, $e=\nabla f/|df|$ at noncritical points, and
$k=(1-t(\sigma)^{3/2})^{1/2}$. Define
\begin{equation}
 \mathcal P_{ee}=2k,\qquad
 \mathcal P|_{e^\perp}=\frac{3+k^2}{2k}I,\qquad
 \mathcal P_{e\perp}=0,
 \qquad L=\mathcal P-\tfrac12(\tr_g\mathcal P)g.
 \label{eq:str-stress-definition}
\end{equation}
At a critical point set $\mathcal P=2g$ and $L=-g$. These fields are
continuous, smooth away from the critical set, and satisfy
\begin{equation}
 \operatorname{div}_g\mathcal P=0,
 \qquad (\tr_gL)^2-|L|_g^2=6,
 \qquad L<0.
 \label{eq:str-stress-identities}
\end{equation}
If $K-(\tr_gK)g\geq0$, then $\mathcal P:K\leq0$. Therefore
$(g,K+L)$ satisfies the flat-cosmological DEC whenever $(g,K)$ satisfies
the $\Lambda=-3$ DEC, with at least the original strict margin, and
its future boundary expansion is strictly smaller.

On any compact set where the DEC and expansion have strict margins,
$\mathcal P$ can be smoothed preserving those margins. This is valid
also at an original boundary face, using only the exterior side.
\end{lemma}

\begin{proof}
Let
\[
 W(\sigma)=\frac{3+k^2}{2k}.
\]
Differentiating $\sigma=\sqrt t/k$ and $k^2=1-t^{3/2}$ gives
\begin{equation}
 W'(\sigma)=\tfrac32t(\sigma),\qquad
 W(\sigma)-\sigma W'(\sigma)=2k.
 \label{eq:str-domain-stress-calculation}
\end{equation}
Consequently $W$ is convex and
$\mathcal P=Wg-\sigma W'e\otimes e$. The weak equation says that
$f$ is a local minimizer of $\int W(|df|)\,dV_g$: convexity gives
the minimizing inequality against every compactly supported additive
variation. Composing with a compactly supported smooth diffeomorphism
is another admissible local variation. Differentiating this inner
variation yields
$\int\mathcal P:\nabla X\,dV_g=0$ for every compactly supported
smooth vector field $X$. The differentiation is valid for the locally
Lipschitz minimizer by change of variables and dominated convergence.
This proves distributional conservation.

Trace reversal in Equation~\eqref{eq:str-stress-definition} gives
\begin{equation}
 L_{ee}=-\frac{3-k^2}{2k},\qquad L|_{e^\perp}=-kI.
 \label{eq:str-L-eigenvalues}
\end{equation}
Both eigenvalues are negative and their quadratic constraint
combination is $6$. As $\sigma\to0$, $k=1+O(\sigma^3)$, so the
definitions extend continuously across critical points, independently
of $e$.

Positivity of $K-(\tr K)g$ implies
$\tr K\leq0$ and $\tr_{e^\perp}K\leq0$. Thus
\[
 \mathcal P:K=2k\tr K+
      \frac{3(1-k^2)}{2k}\tr_{e^\perp}K\leq0.
\]
The flat Hamiltonian of $K+L$ is
\[
 R_g+(\tr(K+L))^2-|K+L|^2
 =\mathcal H_g(K)-2\mathcal P:K,
\]
and its momentum is $\mathcal J_g(K)$ by conservation. Negative
definiteness of $L$ gives the strict boundary expansion decrease.

We spell out the smoothing assertion because pointwise tensor
approximation alone would not control momentum. In a coordinate
chart let $Q^{ij}=\sqrt{\det g}\,\mathcal P^{ij}$. Conservation means
\[
 \partial_jQ^{ij}=-\Gamma^i_{jk}Q^{jk}.
\]
Convolution gives smooth symmetric densities converging uniformly to
$Q$, whose divergence differs from the required connection term by a
commutator tending uniformly to zero. In a boundary chart $x^3\geq0$
use instead
\[
 Q_\epsilon(x)=\int\rho_\epsilon(y)
                         Q(x+2\epsilon e_3-y)\,dy.
\]
Its kernel lies on the original side even at the face, and the same
commutator estimate holds. After a partition of unity, additional
divergence errors are sums of $d\chi\,(Q_\epsilon-Q)$, which also
tend uniformly to zero. Divide by the original volume density and
trace reverse. The algebraic terms and the momentum therefore converge
uniformly on the compact set. Strict DEC and strict negative expansion
survive sufficiently fine smoothing. A separate smoothing scale may
be chosen for each member of a sequence.
\end{proof}

\subsection{Turning the auxiliary end into an asymptotically flat end}
\label{str-bc-section}

We give the end construction with its mass and area estimates. All spherical
derivatives in this subsection use the unit round metric $\gamma$. A dot denotes
$\partial_{\log s}$. The notation $O_j(s^{-a})$ includes $j$ derivatives under
$s\partial_s$ and spherical differentiation, measured in the indicated
orthonormal frame. Constants in estimates for a completed construction may
depend on its fixed parameters.

\begin{lemma}[Constraints in a spherical foliation]
\label{str-bc-constraints}
\label{lem:str-bend-constraints}
Let
\[
 q=F^{-1}ds^2+s^2\gamma,\qquad F>0,\qquad n=\sqrt F\,\partial_s.
\]
Write the orthonormal blocks of a symmetric tensor $B$ as
\[
 B_{nn}=-h,\qquad B_{nT}=G,\qquad B_T=-kI+U,
 \qquad \tr_\gamma U=0.
\]
Here $G$ and $U$ are represented by a one-form and a symmetric tracefree
two-tensor on the unit sphere. Put
$\mathcal H=R_q+(\tr_qB)^2-|B|_q^2$ and
$\mathcal D=\divg_q(B-(\tr_qB)q)$. Then
\begin{equation}
 R_q=\frac{2(1-F)}{s^2}-\frac{2\partial_sF}{s}
       -\frac{2\sqrt F}{s^2}\Delta_\gamma F^{-1/2}.
 \label{str-bc-scalar}
\end{equation}
In particular, if $F=1+s^2k^2-2m/s$, then
\begin{equation}
 \mathcal H=\frac{4\partial_sm}{s^2}
 -\frac{2\sqrt F}{s^2}\Delta_\gamma F^{-1/2}
 +4k(h-k-s\partial_sk)-2|G|^2-|U|^2.
 \label{str-bc-hamiltonian}
\end{equation}
For arbitrary $F$ the momentum components are
\begin{align}
 \mathcal D_n&=\frac{2\sqrt F}{s}(k+s\partial_sk-h)
 +\frac{\divg_\gamma G-2G\cdot\nabla_\gamma\log\sqrt F}{s},
 \label{str-bc-momentum-normal}\\
 \mathcal D_T&=\sqrt F(\partial_sG+3G/s)
 +\frac{1}{s}\bigl\{\nabla_\gamma(h+k)
 -(h-k)\nabla_\gamma\log\sqrt F\bigr\}\nonumber\\
 &\hspace{12mm}+\frac{1}{s}\bigl\{\divg_\gamma U
                    -U\nabla_\gamma\log\sqrt F\bigr\}.
 \label{str-bc-momentum-tangent}
\end{align}
\end{lemma}

\begin{proof}
Choose a spherical orthonormal frame $\bar E_a$, independent of $s$, and set
$E_a=s^{-1}\bar E_a$. At a point where the spherical connection vanishes,
the connection of $q$ is
\[
 \begin{aligned}
 \nabla_{E_a}n&=\frac{\sqrt F}{s}E_a,&
 \nabla_{E_a}E_b&=-\frac{\sqrt F}{s}\delta_{ab}n,\\
 \nabla_nn&=\frac1s(\nabla_\gamma\log\sqrt F)^aE_a,&
 \nabla_nE_a&=-\frac1s(\partial_a\log\sqrt F)n.
 \end{aligned}
\]
Away from this point the second identity also contains the spherical
connection divided by $s$. Thus the leaf shape operator is
$\sqrt F I/s$. For lapse $N=F^{-1/2}$ the scalar Gauss identity is
\[
 R_q=R_{s^2\gamma}-|A|^2-H^2-2n(H)-2N^{-1}\Delta_{s^2\gamma}N.
\]
Substituting $H=2\sqrt F/s$ gives
Equation~\eqref{str-bc-scalar}. Moreover,
\[
 (\tr B)^2-|B|^2=4hk+2k^2-2|G|^2-|U|^2,
\]
which proves Equation~\eqref{str-bc-hamiltonian} by differentiating $F$.
For $P=B-(\tr B)q$, its blocks are $P_{nn}=2k$, $P_{nT}=G$ and
$P_T=(h+k)I+U$. In the normal divergence, the two acceleration terms
give $-2G\cdot\nabla_\gamma\log\sqrt F/s$, while the shape terms give
$2\sqrt F(k-h)/s$. In the tangential divergence the shape terms give
$3\sqrt F G/s$ and the acceleration terms give
$((k-h)I-U)\nabla_\gamma\log\sqrt F/s$. Adding the ordinary block
derivatives proves the two momentum identities.
\end{proof}

\begin{proposition}[Bending with arbitrarily small mass cost]
\label{str-bc-bending}
\label{prop:str-af-bending}
Suppose $(\Omega,g,B)$ is smooth, has one end, has strictly negative future
boundary expansion, and satisfies
$\mathcal H_g-2|\mathcal D_g|_g>0$ everywhere. Suppose its end is
$g=e^{4v}b$, where
\begin{equation}
 v=r^{-3}v_3(\omega)-d_0r^{-3-\delta}+O_3(r^{-4}),
 \qquad 0<\delta<1,\quad d_0>0,
 \label{str-bc-prepared-v}
\end{equation}
with smooth angular coefficients. Suppose also that on this end there are a
positive smooth $C(\omega)$ and a unit covector $e$ such that
\begin{equation}
 B=-g+\frac{C^3}{r^3}
       \left(\frac12g-\frac32e\otimes e\right)+O_1(r^{-4}),
 \qquad e\otimes e=n_r^\flat\otimes n_r^\flat+O_1(r^{-1}),
 \label{str-bc-prepared-B}
\end{equation}
where $n_r$ is the outward unit normal to the coordinate spheres.
The scalar remainder has three symbol derivatives so that the change to the
area coordinate has two controlled metric derivatives. The initial
interpolation needs two derivatives of the metric error and one of the
tensor error; the preparation in Lemma~\ref{lem:str-preparation} supplies
the stronger scalar expansion with all symbol derivatives.

For every $\varepsilon>0$ and every fixed compact subset there exist smooth
data $(q_\varepsilon,B_\varepsilon)$ agreeing with $(g,B)$ on that subset
and near the original boundary, such that
\begin{align}
 &q_\varepsilon\geq(1-\varepsilon)g,
 \label{str-bc-lower-comparison}\\
 &R_{q_\varepsilon}+(\tr B_\varepsilon)^2-|B_\varepsilon|^2
 -2|\divg(B_\varepsilon-(\tr B_\varepsilon)q_\varepsilon)|
 \geq c_\varepsilon(1+s)^{-3-\delta}>0,
 \label{str-bc-strict-dec}\\
 &q_\varepsilon-\delta_{\mathrm{Eucl}}=O_\infty(s^{-1}),\qquad
 R_{q_\varepsilon}=O_\infty(s^{-3-\delta}),\nonumber\\
 &B_\varepsilon=O_\infty(s^{-2}),\qquad
 \tr_{q_\varepsilon}B_\varepsilon=0
 \quad\hbox{on the final end}.
 \label{str-bc-af-bounds}
\end{align}
Its metric ADM energy satisfies
\begin{equation}
 \fint_{\mathbb S^2}(8v_3-C^3/2)
 \leq E(q_\varepsilon)
 \leq\fint_{\mathbb S^2}(8v_3-C^3/2)+\varepsilon.
 \label{str-bc-energy}
\end{equation}
Consequently its full-cut area infimum is at least
$(1-\varepsilon)A_g(S)$.
\end{proposition}

\begin{proof}
We may assume $0<\varepsilon<1$. All modifications take place beyond a radius
chosen at the end of the proof.

\paragraph{Area coordinate and initial coefficients.}
Set $s=e^{2v}r$. Its radial derivative is positive far out. In coordinates
$(s,\omega)$ the metric is
\[
 g=F_g^{-1}(ds-2s\,d_\omega v)^2+s^2\gamma,
 \qquad F_g=(1+r^2)(1+2r\partial_rv)^2,
\]
where the angular differential on the right is initially taken at fixed $r$.
Expansion of Equation~\eqref{str-bc-prepared-v} gives
\begin{equation}
 F_g=1+s^2-\frac{16v_3}{s}
       +4(4+\delta)d_0s^{-1-\delta}+O_2(s^{-2}).
 \label{str-bc-area-coefficient}
\end{equation}
The mixed term has relative orthonormal size
$|d_\omega v|/\sqrt{F_g}=O_2(s^{-4})$. Define
\begin{equation}
 \begin{gathered}
 m_0=8v_3,\qquad d=(8+2\delta)d_0,\qquad
 \mu=m_0-C^3/2,\\
 k=1-\frac{C^3}{2s^3},\quad h=k+s\partial_sk,
 \quad m=\mu-ds^{-\delta}.
 \end{gathered}
 \label{str-bc-initial-coefficients}
\end{equation}
Then $F=1+s^2k^2-2m/s$ agrees with
Equation~\eqref{str-bc-area-coefficient} up to an $O_2(s^{-2})$ term.
Thus $F^{-1}ds^2+s^2\gamma$ differs from $g$ by $O_2(s^{-4})$ in
hyperbolic tensor norm. Equation~\eqref{str-bc-prepared-B} gives the blocks
$B_T=-kI+O_1(s^{-4})$, $B_{nn}=-h+O_1(s^{-4})$,
and $B_{nT}=O_1(s^{-4})$. In particular the approximate spherical pair
differs from the original pair by errors of exactly the orders needed by
the constraints.

For completeness, the time mass of $e^{4v}b$ in the stipulated hyperbolic
metric flux is $\fint 8v_3$. Indeed, for $e=\lambda b$ its flux integrand
is $-2V_0d\lambda+2\lambda dV_0$, and
$\lambda=e^{4v}-1=4v+O(r^{-6})$. The leading integrand on $r=R$ is
$32v_3R^{-2}+o(R^{-2})$, whose integral divided by $16\pi$ is
$8\fint v_3$. This uses precisely the metric time flux; no charge of $B$
is included.

\paragraph{Exactification and finite angular approximation.}
Over a fixed interval in $\log s$, interpolate the metric and tensor to the
pair in Equation~\eqref{str-bc-initial-coefficients}. For example interpolate
the metric covariantly, and interpolate $B+q$ in the same fixed coordinates.
The interpolation is positive definite for a sufficiently distant annulus.
Its constraint error is $O(s^{-4})$, since physical derivatives of the
logarithmic cutoffs are bounded on this hyperbolic part. In the spherical
pair the mass term supplies
\begin{equation}
 \frac{4\partial_s(-ds^{-\delta})}{s^2}
       =4\delta d\,s^{-3-\delta}.
 \label{str-bc-margin}
\end{equation}
All remaining terms in its constraints are $O(s^{-4})$: in particular
$\nabla_\gamma k=O(s^{-3})$, $h-k=O(s^{-3})$,
and $\nabla_\gamma F=O(s^{-1})$. The margin in
Equation~\eqref{str-bc-margin} therefore absorbs the interpolation error.
Next cut $k$ to $1$ on another fixed logarithmic annulus, retaining
$h=k+s\partial_sk$, $G=U=0$, and the same $m$. The new terms are again
$O(s^{-4})$ and are absorbed by the same margin.

Fix $\eta>0$. Choose a finite spherical harmonic sum $\mu_N$, with the same
mean as $\mu$, such that $\|\mu_N-\mu\|_{C^4}<\eta$. Such sums approximate
a smooth function in $C^4$: repeated integration by parts with $\Delta_\gamma$
makes its harmonic coefficients decay faster than every power of the degree,
and hence the differentiated series converges uniformly. On one further
logarithmic annulus replace $\mu$ by
$\mu+\chi(\mu_N-\mu)$, for a nondecreasing cutoff $\chi$. Add to $m$ a
radial nondecreasing function $\rho$ with
$\dot\rho=C\eta\dot\chi$. The signed contribution
$4\dot\chi(\mu_N-\mu)$ to $s^3\mathcal H$ is bounded by
$4\eta\dot\chi$, so a fixed $C>1$ absorbs it. Here $k=1$ and $t=sk=s$;
the remaining scaled angular errors are $O(s^{-2})$. This step costs at
most $C\eta$ in mean mass. From this point onward the freely prescribed
angular coefficient fields belong to one fixed finite collection of
harmonic spaces. Nonlinear expressions such as $F^{-1}$ need not be
finite harmonic sums; their angular derivatives are estimated directly
in the formulas below.

\paragraph{Reduction from large radial boost.}
For radial $t=sk$, put
\[
 F=1+t^2-2m/s,\qquad h=\dot t/s,\qquad G=U=0,
 \qquad m=\mu_N-ds^{-\delta}+\rho(s).
\]
The constraint identities become the exact formulas
\begin{equation}
 \begin{split}
 s^3\mathcal H&=4\dot m-\frac{2\Delta_\gamma m}{F}
           -\frac{6|\nabla_\gamma m|^2}{sF^2},\\
 \mathcal D_n&=0,\qquad
 s^3\mathcal D_T=\frac{(\dot t-t)\nabla_\gamma m}{F}.
 \end{split}
 \label{str-bc-large-exact}
\end{equation}
For example these follow from
$\sqrt F\Delta_\gamma F^{-1/2}
=(\Delta_\gamma m)/(sF)+3|\nabla_\gamma m|^2/(s^2F^2)$.
Thus no expansion involving an unbounded $\dot t$ is needed.

Choose a large fixed $P$. Over a bounded logarithmic interval turn $t=s$
into a constant of comparable size, then lower that constant monotonically
to $P$ over a second bounded logarithmic interval. Explicitly, starting at
$s=R_1$, take $x=\log(s/R_1)$ and
$t=R_1\exp(\int_0^x\kappa(u)\,du)$, where $\kappa$ decreases smoothly
from one to zero and is constant near the endpoints. Its final value is
$T\asymp R_1$. On the next interval take
$t=P+(T-P)(1-\chi(x))$ with a nondecreasing cutoff flat at its endpoints.
These profiles have $P\leq t\leq s$. In the first interval
$t$ is comparable with its starting radius and $|\dot t|\leq Ct$; in the
second interval $t$ is nonincreasing. Accordingly
\begin{equation}
 \int\frac{1+t+|\dot t|}{1+t^2}\,d\log s\leq\frac{C}{P}.
 \label{str-bc-large-cost}
\end{equation}
The contribution of $1+t$ follows from the bounded interval lengths; for
the monotone part the derivative contribution is
$\int_P^T(1+t^2)^{-1}dt\leq P^{-1}$.

As long as $|m|$ is bounded and the starting radius is large,
$F\geq(1+t^2)/2$. Equations~\eqref{str-bc-large-exact} show that a smooth
nonnegative radial addition with
\begin{equation}
 \dot\rho\geq C_N\frac{1+t+|\dot t|}{1+t^2}
 \label{str-bc-large-repair}
\end{equation}
where the angular errors need repair makes their contribution to
$\mathcal H-2|\mathcal D|$ nonnegative. A smooth majorant of $|\dot t|$
can be used, with the same integral bound. The addition starts before $t$
departs from $s$; in this initial overlap the errors already have the smaller
orders treated above. Its cost is at most $C_N/P$. This also verifies the
boundedness of $m$ used in the argument: choose $P$ so that this cost is at
most one and then take the radius large. Continue this repair into a short
overlap on which $t=P$.

\paragraph{The bounded radial boost equations.}
Now $t$ follows a fixed smooth profile from $P$ to zero on a bounded
logarithmic interval. Write $S_0=\sqrt{1+t^2}$ and allow
\begin{equation}
 G=b/s^2,\qquad U=u/s^2,\qquad
 h=\dot t/s+\frac{\divg_\gamma b}{2S_0s^2},
 \label{str-bc-bounded-ansatz}
\end{equation}
where $u$ is tracefree. Decompose $m=\bar\mu+\rho-ds^{-\delta}+m_{\rm ang}$,
with $m_{\rm ang}$ of zero mean. For bounded coefficients and their needed
derivatives, $F=S_0^2+O(s^{-1})$. Substitution in
Lemma~\ref{str-bc-constraints} gives
\begin{align}
 s^3\mathcal H&=4\dot\rho+4\delta ds^{-\delta}
 +4\dot m_{\rm ang}-\frac{2\Delta_\gamma m}{S_0^2}
 +\frac{2t\divg_\gamma b}{S_0}+O(s^{-1}),
 \label{str-bc-bounded-H}\\
 s^3\mathcal D_n&=O(s^{-1}),\nonumber\\
 s^3\mathcal D_T&=S_0(\dot b+b)
 +\frac{\nabla_\gamma\divg_\gamma b}{2S_0}
 +\frac{(\dot t-t)\nabla_\gamma m}{S_0^2}
 +\divg_\gamma u+O(s^{-1}).
 \label{str-bc-bounded-D}
\end{align}
For instance the two leading terms in the normal component are
$-\sqrt F\divg b/(S_0s^3)$ and $\divg b/s^3$, which cancel to the
stated order. The term $4k(h-k-s\partial_sk)$ gives
$2t\divg b/(S_0s^3)$. In the tangential component,
$\sqrt F(\partial_sG+3G/s)=\sqrt F(\dot b+b)/s^3$,
and $\nabla\log\sqrt F=-\nabla m/(sF)$. These observations account for
every leading term in Equations~\eqref{str-bc-bounded-H} and
\eqref{str-bc-bounded-D}; the quadratic $G,U$ terms are $O(s^{-4})$.

We cancel the leading angular terms by solving
\begin{align}
 \dot m_{\rm ang}&=\frac{\Delta_\gamma m}{2S_0^2}
                   -\frac{t\divg_\gamma b}{2S_0},
 \label{str-bc-angular-mass}\\
 S_0(\dot b+b)&+\frac{\nabla_\gamma\divg_\gamma b}{2S_0}
 +\frac{(\dot t-t)\nabla_\gamma m}{S_0^2}
 +\divg_\gamma u=0.
 \label{str-bc-angular-momentum}
\end{align}
Let $Y$ be a spherical harmonic with
$\Delta_\gamma Y=-\lambda_\ell Y$,
$\lambda_\ell=\ell(\ell+1)$. Commutation of a covariant derivative with
the Hessian, using $\Ric_\gamma=\gamma$, gives
\begin{equation}
 \divg_\gamma\left(\Hess_\gamma Y+
                   \frac{\lambda_\ell}{2}Y\gamma\right)
       =\left(1-\frac{\lambda_\ell}{2}\right)\nabla_\gamma Y.
 \label{str-bc-tf-hessian}
\end{equation}
For $\ell\geq2$ set the corresponding component of $b$ equal to zero.
Its mass coefficient $M_\ell$ and tracefree tensor are then
\begin{equation}
 \dot M_\ell=-\frac{\lambda_\ell M_\ell}{2S_0^2},\qquad
 u_\ell=-\frac{(\dot t-t)M_\ell}
 {S_0^2(1-\lambda_\ell/2)}
 \left(\Hess_\gamma Y+\frac{\lambda_\ell}{2}Y\gamma\right).
 \label{str-bc-higher-modes}
\end{equation}
For each dipole write $m_{\rm ang}=MY$ and $b=A\nabla_\gamma Y$.
Since $\lambda_1=2$, the remaining equations are
\begin{equation}
 \dot M=-\frac{M}{S_0^2}+\frac{tA}{S_0},\qquad
 \dot A=-\frac{t^2A}{S_0^2}+\frac{(t-\dot t)M}{S_0^3}.
 \label{str-bc-dipole-ode}
\end{equation}
This finite linear system has a unique smooth solution throughout the
chosen bounded interval. All of its norms are finite once the profile,
$P$, and the harmonic cutoff have been fixed. Its equations preserve the
mean of $m_{\rm ang}$.

To join smoothly to the previous stage, use an overlap of fixed logarithmic
length on which $t=P$. Start with the prescribed mass coefficients and
$A=0$, and multiply the right sides of the derivative equations
\eqref{str-bc-higher-modes} and \eqref{str-bc-dipole-ode}, and the formula
for $u_\ell$, by a cutoff increasing smoothly from zero to one. At constant
$P$, the equations imply on this overlap
\[
 |A|+|\dot A|\leq C_NP^{-2},\qquad
 |\dot M_\ell|\leq C_NP^{-2},\qquad |u_\ell|\leq C_NP^{-1}.
\]
The first estimate follows directly from the scalar equation for $A$ by
variation of constants; its source is $PM/S_0^3=O(P^{-2})$ and its
zeroth-order coefficient is bounded. The mass coefficients remain bounded
by the same finite-dimensional estimate. Inserting these bounds into the
uncancelled parts of Equations~\eqref{str-bc-bounded-H} and
\eqref{str-bc-bounded-D} bounds them by $C_N/P$. Continue the radial repair
with a smooth derivative of this size until the cutoff is identically one.
It can decrease to zero in proportion to one minus that cutoff; its total
cost is at most $C_N/P$. Thereafter the displayed angular equations hold
exactly and no further radial mass bump is needed.

\paragraph{The final dipole adjustment.}
Take $t=0$ after the preceding profile, flat at the join. Equations
\eqref{str-bc-higher-modes} and \eqref{str-bc-dipole-ode} then give
$\dot M_\ell=-\lambda_\ell M_\ell/2$, $\dot A=0$, and $u=0$.
Let $C_1$ denote the sum of the resulting dipole potentials, so initially
$b=\nabla_\gamma C_1$. On one further bounded logarithmic interval set
\begin{equation}
 F=1-2m/s,\quad k=a/s^2,\quad G=b/s^2,\quad U=0,\quad
 h=(\dot a-a+\divg_\gamma b/2)/s^2,
 \label{str-bc-final-ansatz}
\end{equation}
where
\begin{equation}
 a=\zeta C_1/2,\qquad b=\nabla_\gamma A_1,
 \qquad A_1=C_1-a,
 \label{str-bc-final-interpolation}
\end{equation}
and $\zeta$ rises smoothly from zero to one, flat at both endpoints.
Continue $\dot m_{\rm ang}=\Delta_\gamma m/2$ and keep $\rho$ constant.
The leading tangential momentum in this convention is
$s^{-3}\nabla_\gamma\partial_{\log s}(a+A_1)$, which vanishes identically.
The normal leading term also cancels. There is no new leading Hamiltonian
term: direct substitution of Equation~\eqref{str-bc-final-ansatz} gives
\begin{equation}
 \mathcal H=\frac{4\dot m-2\Delta_\gamma m/F}{s^3}
 +\frac{-6|\nabla_\gamma m|^2/F^2+4a\dot a-2a^2
       +2a\divg_\gamma b-2|b|^2}{s^4}.
 \label{str-bc-final-H}
\end{equation}
In particular all new errors are $O(s^{-4})$. The formula uses
$F=1-2m/s$, not $1+s^2k^2-2m/s$; the difference belongs to the displayed
quadratic order. The join is smooth because initially $a=0$, and the
preceding $t$ profile is already zero with all derivatives.

At the endpoint $a=A_1=C_1/2$ is independent of $s$ and
$\Delta_\gamma a=-2a$. Thus
\begin{equation}
 B_{nn}=2a/s^2,\qquad B_T=-aI/s^2,\qquad
 B_{nT}=\nabla_\gamma a/s^2,\qquad \tr_qB=0.
 \label{str-bc-tracefree-tail}
\end{equation}
This is exact tracefreeness, rather than an asymptotic trace estimate.

\paragraph{Error bounds and the order of choices.}
Here is the constraint ledger; its errors are expressed after multiplication
by $s^3$. The positive contribution retained at every stage is
$4\delta d s^{-\delta}$.
\begin{center}
\begin{tabular}{lll}
\toprule
Stage & Error or repair size & Total added mean mass\\
\midrule
Initial interpolation and cutoff of $k-1$ & $C/s$ & $0$\\
Finite harmonic interpolation & $C\eta\dot\chi+C/s$ & $\leq C\eta$\\
Large boost reduction & $C_N(1+t+|\dot t|)/(1+t^2)$ & $\leq C_N/P$\\
Overlap for angular equations & $C_N/P+C_{N,P}/s$ & $\leq C_N/P$\\
Bounded boost and dipole adjustment & $C_{N,P}/s$ & $0$\\
Final infinite tail & $C_{N,P}/s$ & $0$\\
\bottomrule
\end{tabular}
\end{center}
The constants in the last three rows also depend on the fixed smooth
profiles. Choose first $\eta$ and the finite harmonic approximation with
$C\eta<\varepsilon/3$. Next choose $P$ so large that the sum of the two
$C_N/P$ costs is less than $2\varepsilon/3$. Then fix the bounded boost,
overlap, and final dipole profiles and solve their finite linear systems.
Only now choose the starting radius $R$ large enough that $F>0$ and
\begin{equation}
 C_{N,P}s^{-1}\leq\delta d s^{-\delta}\qquad(s\geq R)
 \label{str-bc-radius-choice}
\end{equation}
for every remaining error constant. This is possible because $\delta<1$.
The large boost formulas are exact, so the possibly large derivatives of
that radius-dependent profile do not enter $C_{N,P}$ in
Equation~\eqref{str-bc-radius-choice}. Smooth repairs may overlap the
successive stages as described above; where a repair is turning on before
$t$ departs from $s$, the errors have the initial smaller order, and where
it turns off the angular cutoff is already becoming identically one.
This proves a lower bound $c s^{-3-\delta}$ for the new end. On the
unchanged compact part the original strict inequality has a positive
minimum. Together these give Equation~\eqref{str-bc-strict-dec}.

\paragraph{Decay, energy, and the metric comparison.}
Let $s_*$ be the beginning of the final tail. Its coefficients have the form
\begin{equation}
 m=m_\infty-ds^{-\delta}
    +\sum_{\ell=1}^{N}\sum_j c_{\ell j}
          (s/s_*)^{-\lambda_\ell/2}Y_{\ell j},
 \qquad m_\infty=\fint\mu+\rho_\infty.
 \label{str-bc-final-mass}
\end{equation}
Every nonconstant mode is $O_\infty(s^{-1})$. Since
$\dot m=\delta ds^{-\delta}+\Delta_\gamma m/2$, the exact scalar formula is
\begin{equation}
 R_q=4\delta ds^{-3-\delta}
 -\frac{4m\Delta_\gamma m}{s^4F}
 -\frac{6|\nabla_\gamma m|^2}{s^4F^2}
 =4\delta ds^{-3-\delta}+O_\infty(s^{-5}).
 \label{str-bc-tail-scalar}
\end{equation}
For example the exact endpoint momentum is
\[
 \begin{split}
 s^3\mathcal D_n&=-2a(1-\sqrt F)
                  +\frac{2\nabla a\cdot\nabla m}{sF},\\
 s^3\mathcal D_T&=(\sqrt F-1)\nabla a
                  -\frac{3a\nabla m}{sF}.
 \end{split}
\]
It is $O_\infty(s^{-4})$. Equations~\eqref{str-bc-final-mass} and
\eqref{str-bc-tracefree-tail} prove all the symbol estimates in
Equation~\eqref{str-bc-af-bounds}; conversion to Cartesian derivatives
loses one power of $s$ per derivative. In particular
$q-\delta_{\mathrm{Eucl}}=(F^{-1}-1)ds^2=O_\infty(s^{-1})$.

For this metric the ADM integral at $s$ equals
$\frac{s}{2}\fint(F^{-1}-1)$. Indeed writing
$q_{ij}=\delta_{ij}+f\omega_i\omega_j$ with $f=F^{-1}-1$, the radial
contraction of $\partial_jq_{ij}-\partial_iq_{jj}$ is $2f/s$;
angular derivatives cancel in this contraction. Its limit is $m_\infty$.
All angular changes have zero mean, and the radial additions are
nonnegative with total at most $\varepsilon$, proving
Equation~\eqref{str-bc-energy}.

Throughout the large boost reduction $t\leq s$, and throughout the bounded
boost $t\leq P\leq s$. Thus $F\leq1+s^2+C/s$ everywhere after the initial
interpolation. On the final dipole adjustment and tail the stronger bound
$F\leq1+C/s$ holds. The original inverse radial coefficient is
$1+s^2+O(s^{-1})$, and its mixed relative term is $O(s^{-4})$.
Consequently comparison of the radial coefficients and the unchanged
$s^2\gamma$ blocks gives $q\geq(1-o_R(1))g$ uniformly on the modified
end. The initial interpolation has the same bound. Increase $R$ once more
to obtain Equation~\eqref{str-bc-lower-comparison}. This lower comparison
also gives completeness. The boundary and the number of ends are unchanged.
Every admissible full cut is the same embedded submanifold before and after
the construction. Restricting $q\geq(1-\varepsilon)g$ to its two-dimensional
tangent planes gives $\Area_q(\Gamma)\geq(1-\varepsilon)\Area_g(\Gamma)$.
Taking the infimum over the full intrinsic boundaries, including $S$ when
$D=\Omega$, proves the asserted area comparison.
\end{proof}

\subsection{Completion of the one-sign inequality}

\begin{proof}[Proof of Theorem~\ref{thm:str-one-sign}]
If $K-(\tr K)g\leq0$, replace $K$ by $-K$. Indeed
$\tr_TK\geq0$ for every two-plane, so this decreases the future
expansion and preserves DEC. We may therefore assume
$K-(\tr K)g\geq0$.

Apply Lemma~\ref{lem:str-preparation} and first fix one prepared data
set, suppressing its index. Add the short inner collar from
Lemma~\ref{lem:str-added-collar}. Solve the saturating Dirichlet problem
there with height $D$, and restrict its stress to the original
exterior. Lemma~\ref{lem:str-stress-identities} gives strict
flat-cosmological DEC and strict future trapping for $(g,K+L)$.
The stress is smooth on an end, and it can be smoothed on the remaining
compact region without changing either the end or the strict
conditions.

Proposition~\ref{prop:str-af-bending}, with arbitrarily small error,
produces a smooth asymptotically Euclidean exterior satisfying all
hypotheses in Equation~\eqref{eq:str-flat-interface}, with
\[
 E\leq p_0(g)-\tfrac12\fint C_D^3+\epsilon,
 \qquad A_q(S)\geq(1-\epsilon)A_g(S).
\]
Theorem~\ref{thm:fl-deformation} therefore gives, on letting
$\epsilon\downarrow0$,
\[
 p_0(g)\geq\sqrt{\frac{A_g(S)}{16\pi}}
                       +\tfrac12\fint C_D^3
 \geq\sqrt{\frac{A_g(S)}{16\pi}}
                       +\tfrac12(\mathcal F_D/(4\pi))^{3/2}.
\]
For each fixed added collar, let $D\to\infty$ and use
Lemma~\ref{lem:str-flux-saturation}. Then let the collar width tend to
zero and use Equation~\eqref{eq:str-collar-area}. This proves
Equation~\eqref{eq:str-component-bound} for the prepared data. Finally,
remove the preparation using the mass and area convergence in
Lemma~\ref{lem:str-preparation}.

Every limiting operation just used is a limit of scalar inequalities.
No smooth limiting stress as $D\to\infty$, or regular limiting metric
at equality, is required for this argument. The fixed-chart proof did
not use a sign or timelikeness assumption on the flux covector. If it
is future timelike, a background hyperbolic isometry makes its time
component equal to its Lorentz norm. The geometric and decay
hypotheses, the compact support and sign of the constraint stress,
and the cut area are unchanged. Lemma~\ref{lem:setup-mass-covariance}
then proves the final assertion.
\end{proof}
